\subsection{DEFINITION OF MATRICES}

DEFINITION 1. Let I, K, H be three sets; a matrix of type (I, K) with elements in H (or a matrix of type (I, K) over H) is any family \(M = (m_{\iota\kappa})_{(\iota, \kappa) \in I \times K}\) of elements of H whose indexing set is the product \(I \times K\) . For all \(\iota \in I\) , the family \((m_{\iota\kappa})_{\kappa \in K}\) is called the row of \(M\) of index \(\iota\) ; for all \(\kappa \in K\) , the family \((m_{\iota\kappa})_{\iota \in I}\) is called the column of \(M\) of index \(\kappa\) .

If I (resp. K) is finite, M is said to be a matrix with a finite number of rows (resp. columns). The set of matrices of type (I, K) over H is identified with the product \(H^{I \times K}\) .

The names ``row'' and ``column'' arise from the fact that, in the case where I and K are intervals \([1,p]\) , \([1,q]\) of N, the elements of the matrix are envisaged as set out in a rectangular array with p rows (arranged horizontally) and q columns (arranged vertically):

\[
\left( \begin{array}{c c c c c} m _ {1 1} & m _ {1 2} & \ldots & m _ {1 q} \\ m _ {2 1} & m _ {2 2} & \ldots & m _ {2 q} \\ \cdot & \cdot & \cdot & \cdot & \cdot \\ m _ {p 1} & m _ {p 2} & \ldots & m _ {p q} \end{array} \right)
\]

When p and q are explicit integers sufficiently small for it to be practicable, it is a convention that the above array is a symbol effectively denoting the matrix in question; this notation enables us to dispense with the use of indices, it being understood that the indices of an element are determined by its place in the array; for example, when we speak of the matrix

\[
\left( \begin{array}{c c c} a & b & c \\ d & e & f \end{array} \right)
\]

we mean the matrix \((m_{ij})_{1\leqslant i\leqslant 2,1\leqslant j\leqslant 3}\) such that

\[
m _ {1 1} = a, m _ {1 2} = b, m _ {1 3} = c, m _ {2 1} = d, m _ {2 2} = e, m _ {2 3} = f.
\]

Instead of matrix of type \(([1, p], [1, q])\) , we also say matrix of type \((p, q)\) , or matrix with \(p\) rows and \(q\) columns, if no confusion arises; the set of matrices of type \((p, q)\) over H is sometimes denoted by \(\mathbf{M}_{p,q}(\mathrm{H})\) .

Every matrix over H for which one of the indexing sets I, K is empty is identical with the empty family of elements of H; it is also called the empty matrix. When \(I = \{i_{0}\}\) (resp. \(K = \{k_{0}\}\) ) is a set consisting of a single element, M is called a row matrix (resp. column matrix) and the row (resp. column) index can then be suppressed in the notation; when I and K are both sets with one element, a matrix of type (I, K) is often identified with the unique element in this matrix.

A subfamily \(M' = (m_{\iota \kappa})_{(\iota, \kappa) \in J \times L}\) of a matrix \(M = (m_{\iota \kappa})_{(\iota, \kappa) \in I \times K}\) , whose indexing set is the product of a subset \(J\) of \(I\) and a subset \(L\) of \(K\) , is called a submatrix of the matrix \(M\) ; it is said to be obtained by suppressing in \(M\) the rows of index \(\iota \notin J\) and the columns of index \(\kappa \notin L\) ; conversely, \(M\) is said to be obtained by bordering \(M'\) with the rows of index \(\iota \notin J\) and the columns of index \(\kappa \notin L\) .

DEFINITION 2. The transpose of a matrix \(M = (m_{\iota \kappa})_{(\iota, \kappa) \in \mathbf{I} \times \mathbf{K}}\) , denoted by \({}^tM\) , is the matrix \((m_{\kappa \iota}')_{(\kappa, \iota) \in \mathbf{K} \times \mathbf{I}}\) over \(\mathbf{H}\) given by \(m_{\kappa \iota}' = m_{\iota \kappa}\) for all \((\iota, \kappa) \in \mathbf{I} \times \mathbf{K}\) .

It follows from this definition that the transpose of a matrix of type (I, K) is a matrix of type (K, I) and that

\[
{ } ^ { t } ( { } ^ { t } M ) = M .\tag{1}
\]

